
\subsubsection*{2.6.3}
\phantomsection
\label{sga3:VIIB:2.6.3}

If \(\mathbf L\) is a flat Lie algebra over \(\mathbf O_k\), the
formal group \(\mathscr G(\mathbf L)\) can be characterized by a
universal property.\footnote{%
\label{sga3:VIIB:nde:111}The editors modified what follows to take
advantage of the addition made in \SGARefLink{sga3:VIIB:2.6.2:prop}{2.6.2}.}
Indeed, every morphism
\[
  \phi:\mathscr G(\mathbf L)\longrightarrow G
\]
to a formal group \(G\) induces a morphism
\[
  \operatorname{Lie}(\phi):
  \operatorname{Lie}\bigl(\mathscr G(\mathbf L)\bigr)
  \longrightarrow
  \operatorname{Lie}(G).
\]
Composing it with
\[
  \tau_{\mathbf L}:
  \mathbf L\longrightarrow
  \operatorname{Lie}\bigl(\mathscr G(\mathbf L)\bigr)
\]
\textup{(cf.\ \SGARefLink{sga3:VIIB:2.6.2:prop}{2.6.2})} gives a morphism
\[
  \phi':\mathbf L\longrightarrow\operatorname{Lie}(G).
\]
We then have:

\medskip
\noindent\textbf{Proposition.}
\phantomsection
\label{sga3:VIIB:2.6.3:prop}
{\itshape Let \(G\) be a formal \(k\)-group and let \(\mathbf L\)
be a flat Lie algebra over \(\mathbf O_k\). The map
\(\phi\mapsto\phi'\) defined above is a bijection
\[
  \operatorname{Hom}_{\operatorname{Grf}/k}
    \bigl(\mathscr G(\mathbf L),G\bigr)
  \xrightarrow{\ \sim\ }
  \operatorname{Hom}_{\operatorname{Lie}}
    \bigl(\mathbf L,\operatorname{Lie}(G)\bigr),
\]
where the right-hand side is the set of morphisms of Lie algebras
over \(\mathbf O_k\) from \(\mathbf L\) to
\(\operatorname{Lie}(G)\).}

\textit{Proof.}
We may immediately reduce to the case in which \(k\) is Artinian.
Put \(L=\mathbf L(k)\). By \SGARefLink{sga3:VIIB:2.3.1}{2.3.1},
\[
  \operatorname{Hom}_{\operatorname{Grf}/k}
    \bigl(\mathscr G(\mathbf L),G\bigr)
\]
is in bijection with the set of unital algebra morphisms
\[
  h:U(L)\longrightarrow\mathbf H(G)
\]
for which the following diagram is commutative:
\begin{center}
  \phantomsection
  \label{sga3:VIIB:2.6.3:universal-property-diagram}
  \begin{tikzcd}[row sep=1.2em,column sep=4.2em]
U(L) \arrow[r,"h"] \arrow[dd,"\Delta_{U(L)}"']
  & \mathbf H(G) \arrow[dr,"\delta_G"] & {} \\
{} & {} & \mathbf H(G\times G) \\
U(L)\otimes U(L) \arrow[r,"h\otimes h"]
  & \mathbf H(G)\otimes\mathbf H(G)
      \arrow[ur,"\varphi_G"']
  & {.}
\end{tikzcd}

\end{center}
Now \(h\) is determined by its restriction to \(L\), which is a
morphism from \(L\) to the Lie algebra underlying
\(\mathbf H(G)\). Commutativity of the diagram says precisely that
\(h\) sends \(L\) into the subset of \(\mathbf H(G)\) formed by
the elements \(x\) satisfying
\[
  \delta_G(x)
  =
  \varphi_G(x\otimes1+1\otimes x).
\]
This subset is exactly \(\operatorname{Lie}(G)\); cf.\ \SGARefLink{sga3:VIIB:2.6:item:c}{2.6(c)}.
The proposition follows.
